\begin{table}[tbp]
\caption{From an observation to a testable advisory question.}
\label{tab:ch10-findings}
\centering\small\setlength{\tabcolsep}{4pt}\renewcommand{\arraystretch}{1.18}
\begin{tabularx}{\linewidth}{@{}>{\raggedright\arraybackslash}p{.22\linewidth}*{3}{>{\raggedright\arraybackslash}X}@{}}
\toprule
Observed finding & Bounded interpretation & Candidate response & Evidence needed next \\
\midrule
Different leading regions across the three maps & Participation, contribution, and project influence describe distinct properties & Investigate the projects and roles behind a regional profile & Aligned periods, population coverage, and attribution histories \\
\addlinespace
Different contribution-origin shares in the two project groups & The reported samples have different source compositions & Examine access to sustained cross-border participation & Comparable projects; contributor entry, retention, and role histories \\
\addlinespace
Hangzhou: 350 to 486, with an intervening decline & Net growth coexists with variation within the period & Investigate which local collaborative functions changed & Individual membership, projects, institutions, and programme exposure \\
\addlinespace
Complementary measures in the Nacos and Mulan dashboards & Project operations and community portfolios require different review scopes & Examine the collaborative functions behind a change & Project histories, role assignments, and maintainer evidence \\
\bottomrule
\end{tabularx}
\tablenote{Findings derive from the report and the supplied Open Source Evaluatology applications. Responses are proposals developed in this chapter, not observed interventions or estimated effects.}
\end{table}
